\documentclass[12pt,a4paper]{article}

\usepackage[a4paper,text={16.5cm,25.2cm},centering]{geometry}
\usepackage{lmodern}
\usepackage{amssymb,amsmath}
\usepackage{bm}
\usepackage{graphicx}
\usepackage{microtype}
\usepackage{hyperref}
\usepackage[usenames,dvipsnames]{xcolor}
\setlength{\parindent}{0pt}
\setlength{\parskip}{1.2ex}




\hypersetup
       {   pdfauthor = {  },
           pdftitle={  },
           colorlinks=TRUE,
           linkcolor=black,
           citecolor=blue,
           urlcolor=blue
       }




\usepackage{upquote}
\usepackage{listings}
\usepackage{xcolor}
\lstset{
    basicstyle=\ttfamily\footnotesize,
    upquote=true,
    breaklines=true,
    breakindent=0pt,
    keepspaces=true,
    showspaces=false,
    columns=fullflexible,
    showtabs=false,
    showstringspaces=false,
    escapeinside={(*@}{@*)},
    extendedchars=true,
}
\newcommand{\HLJLt}[1]{#1}
\newcommand{\HLJLw}[1]{#1}
\newcommand{\HLJLe}[1]{#1}
\newcommand{\HLJLeB}[1]{#1}
\newcommand{\HLJLo}[1]{#1}
\newcommand{\HLJLk}[1]{\textcolor[RGB]{148,91,176}{\textbf{#1}}}
\newcommand{\HLJLkc}[1]{\textcolor[RGB]{59,151,46}{\textit{#1}}}
\newcommand{\HLJLkd}[1]{\textcolor[RGB]{214,102,97}{\textit{#1}}}
\newcommand{\HLJLkn}[1]{\textcolor[RGB]{148,91,176}{\textbf{#1}}}
\newcommand{\HLJLkp}[1]{\textcolor[RGB]{148,91,176}{\textbf{#1}}}
\newcommand{\HLJLkr}[1]{\textcolor[RGB]{148,91,176}{\textbf{#1}}}
\newcommand{\HLJLkt}[1]{\textcolor[RGB]{148,91,176}{\textbf{#1}}}
\newcommand{\HLJLn}[1]{#1}
\newcommand{\HLJLna}[1]{#1}
\newcommand{\HLJLnb}[1]{#1}
\newcommand{\HLJLnbp}[1]{#1}
\newcommand{\HLJLnc}[1]{#1}
\newcommand{\HLJLncB}[1]{#1}
\newcommand{\HLJLnd}[1]{\textcolor[RGB]{214,102,97}{#1}}
\newcommand{\HLJLne}[1]{#1}
\newcommand{\HLJLneB}[1]{#1}
\newcommand{\HLJLnf}[1]{\textcolor[RGB]{66,102,213}{#1}}
\newcommand{\HLJLnfm}[1]{\textcolor[RGB]{66,102,213}{#1}}
\newcommand{\HLJLnp}[1]{#1}
\newcommand{\HLJLnl}[1]{#1}
\newcommand{\HLJLnn}[1]{#1}
\newcommand{\HLJLno}[1]{#1}
\newcommand{\HLJLnt}[1]{#1}
\newcommand{\HLJLnv}[1]{#1}
\newcommand{\HLJLnvc}[1]{#1}
\newcommand{\HLJLnvg}[1]{#1}
\newcommand{\HLJLnvi}[1]{#1}
\newcommand{\HLJLnvm}[1]{#1}
\newcommand{\HLJLl}[1]{#1}
\newcommand{\HLJLld}[1]{\textcolor[RGB]{148,91,176}{\textit{#1}}}
\newcommand{\HLJLs}[1]{\textcolor[RGB]{201,61,57}{#1}}
\newcommand{\HLJLsa}[1]{\textcolor[RGB]{201,61,57}{#1}}
\newcommand{\HLJLsb}[1]{\textcolor[RGB]{201,61,57}{#1}}
\newcommand{\HLJLsc}[1]{\textcolor[RGB]{201,61,57}{#1}}
\newcommand{\HLJLsd}[1]{\textcolor[RGB]{201,61,57}{#1}}
\newcommand{\HLJLsdB}[1]{\textcolor[RGB]{201,61,57}{#1}}
\newcommand{\HLJLsdC}[1]{\textcolor[RGB]{201,61,57}{#1}}
\newcommand{\HLJLse}[1]{\textcolor[RGB]{59,151,46}{#1}}
\newcommand{\HLJLsh}[1]{\textcolor[RGB]{201,61,57}{#1}}
\newcommand{\HLJLsi}[1]{#1}
\newcommand{\HLJLso}[1]{\textcolor[RGB]{201,61,57}{#1}}
\newcommand{\HLJLsr}[1]{\textcolor[RGB]{201,61,57}{#1}}
\newcommand{\HLJLss}[1]{\textcolor[RGB]{201,61,57}{#1}}
\newcommand{\HLJLssB}[1]{\textcolor[RGB]{201,61,57}{#1}}
\newcommand{\HLJLnB}[1]{\textcolor[RGB]{59,151,46}{#1}}
\newcommand{\HLJLnbB}[1]{\textcolor[RGB]{59,151,46}{#1}}
\newcommand{\HLJLnfB}[1]{\textcolor[RGB]{59,151,46}{#1}}
\newcommand{\HLJLnh}[1]{\textcolor[RGB]{59,151,46}{#1}}
\newcommand{\HLJLni}[1]{\textcolor[RGB]{59,151,46}{#1}}
\newcommand{\HLJLnil}[1]{\textcolor[RGB]{59,151,46}{#1}}
\newcommand{\HLJLnoB}[1]{\textcolor[RGB]{59,151,46}{#1}}
\newcommand{\HLJLoB}[1]{\textcolor[RGB]{102,102,102}{\textbf{#1}}}
\newcommand{\HLJLow}[1]{\textcolor[RGB]{102,102,102}{\textbf{#1}}}
\newcommand{\HLJLp}[1]{#1}
\newcommand{\HLJLc}[1]{\textcolor[RGB]{153,153,119}{\textit{#1}}}
\newcommand{\HLJLch}[1]{\textcolor[RGB]{153,153,119}{\textit{#1}}}
\newcommand{\HLJLcm}[1]{\textcolor[RGB]{153,153,119}{\textit{#1}}}
\newcommand{\HLJLcp}[1]{\textcolor[RGB]{153,153,119}{\textit{#1}}}
\newcommand{\HLJLcpB}[1]{\textcolor[RGB]{153,153,119}{\textit{#1}}}
\newcommand{\HLJLcs}[1]{\textcolor[RGB]{153,153,119}{\textit{#1}}}
\newcommand{\HLJLcsB}[1]{\textcolor[RGB]{153,153,119}{\textit{#1}}}
\newcommand{\HLJLg}[1]{#1}
\newcommand{\HLJLgd}[1]{#1}
\newcommand{\HLJLge}[1]{#1}
\newcommand{\HLJLgeB}[1]{#1}
\newcommand{\HLJLgh}[1]{#1}
\newcommand{\HLJLgi}[1]{#1}
\newcommand{\HLJLgo}[1]{#1}
\newcommand{\HLJLgp}[1]{#1}
\newcommand{\HLJLgs}[1]{#1}
\newcommand{\HLJLgsB}[1]{#1}
\newcommand{\HLJLgt}[1]{#1}


\def\endash{–}
\def\bbD{ {\mathbb D} }
\def\bbZ{ {\mathbb Z} }
\def\bbR{ {\mathbb R} }
\def\bbC{ {\mathbb C} }

\def\x{ {\vc x} }
\def\a{ {\vc a} }
\def\b{ {\vc b} }
\def\c{ {\vc c} }
\def\e{ {\vc e} }
\def\k{ {\vc k} }
\def\f{ {\vc f} }
\def\g{ {\vc g} }
\def\h{ {\vc h} }
\def\u{ {\vc u} }
\def\v{ {\vc v} }
\def\y{ {\vc y} }
\def\z{ {\vc z} }
\def\w{ {\vc w} }

\def\bt{ {\tilde b} }
\def\ct{ {\tilde c} }
\def\Ut{ {\tilde U} }
\def\Qt{ {\tilde Q} }
\def\Rt{ {\tilde R} }
\def\Xt{ {\tilde X} }
\def\acos{ {\rm acos}\, }
\def\sinc{ {\rm sinc}\, }

\def\red#1{ {\color{red} #1} }
\def\blue#1{ {\color{blue} #1} }
\def\green#1{ {\color{ForestGreen} #1} }
\def\magenta#1{ {\color{magenta} #1} }

\input{somacros}

\begin{document}



\textbf{Numerical Analysis MATH50003 (2025\ensuremath{\endash}26) Problem Sheet 9}

We explore some properties of discrete convolution, its relationship to continuous convolution of periodic functions, and cross-correlation. We also investigate the Chebyshev U polynomials, which are closely related to Chebyshev T polynomials.

We saw in the notes that discrete convolutions can be diagonalised in terms of the Discrete Fourier Transform (DFT). Here we use this property to prove two basic results:

\rule{\textwidth}{1pt}
\textbf{Problem 1} Use the connection with the DFT to prove the following properties of convolutions for $\ensuremath{\bm{\f}},\ensuremath{\bm{\g}},\ensuremath{\bm{\h}} \ensuremath{\in} \ensuremath{\bbC}^n$:

\begin{itemize}
\item[1. ] \emph{Commutativity}: $\ensuremath{\bm{\f}} \ensuremath{\star} \ensuremath{\bm{\g}} = \ensuremath{\bm{\g}} \ensuremath{\star} \ensuremath{\bm{\f}}$.


\item[2. ] \emph{Associativity}: $(\ensuremath{\bm{\f}} \ensuremath{\star} \ensuremath{\bm{\g}}) \ensuremath{\star} \ensuremath{\bm{\h}} = \ensuremath{\bm{\f}} \ensuremath{\star} (\ensuremath{\bm{\g}} \ensuremath{\star} \ensuremath{\bm{\h}})$.

\end{itemize}
\rule{\textwidth}{1pt}
For 2\ensuremath{\pi}-periodic functions $a,f$ we have a continuous periodic convolution given by
\[
a \ensuremath{\star} f := {1 \over 2\ensuremath{\pi}} \ensuremath{\int}_0^{2\ensuremath{\pi}}  a(\ensuremath{\theta}-\ensuremath{\varphi}) f(\ensuremath{\varphi}) \D \ensuremath{\varphi}.
\]
Note the standard continuous convolution involves an integral over $\ensuremath{\bbR}$ but because the functions do not decay that no longer makes sense. A natural way to discretise this integral is to use the Trapezium rule. Here we see how this can be related to the discrete convolution and the error understood in terms of the true Fourier coefficients.

\rule{\textwidth}{1pt}
\textbf{Problem 2}  For $2\ensuremath{\pi}$-periodic smooth functions $a$ and $f$ (i.e., such that $a(\ensuremath{\theta}+2\ensuremath{\pi}) = a(\ensuremath{\theta})$ and $b(\ensuremath{\theta}+2\ensuremath{\pi}) = b(\ensuremath{\theta})$) with no negative Fourier coefficients ($0 = \hat{a}_{-1} = \hat{a}_{-2} = \ensuremath{\ldots}$ and $0 = \hat{f}_{-1} = \hat{f}_{-2} = \ensuremath{\ldots}$ ) consider the periodic Trapezium rule approximation to the (continuous) convolution:
\[
\underbrace{ {1 \over 2\ensuremath{\pi}} \ensuremath{\int}_0^{2\ensuremath{\pi}}  a(\ensuremath{\theta}-\ensuremath{\varphi}) f(\ensuremath{\varphi}) \D \ensuremath{\varphi}}_{=: (a \ensuremath{\star} f)(\ensuremath{\theta})} \ensuremath{\approx} \underbrace{\ensuremath{\sum}_{j=0}^{n-1} a(\ensuremath{\theta}-\ensuremath{\theta}_j) f(\ensuremath{\theta}_j)}_{
=: (a \ensuremath{\star}_n f)(\ensuremath{\theta})}
\]
for $\ensuremath{\theta}_j = 2\ensuremath{\pi}j/n$. Express the Fourier coefficients of $a \ensuremath{\star}_n f$  in terms of $\hat{a}_k$ and $\hat{f}_k$ and thereby prove uniform convergence as $n \ensuremath{\rightarrow} \ensuremath{\infty}$ provided both $a$ and $f$ have absolutely summable Fourier coefficients. Hint: Integrals and infinite sums can be swapped due to uniform and absolute convergence.

\rule{\textwidth}{1pt}
Convolutional Neural Networks (CNNs) are typically defined in terms of \emph{cross-correlation} instead of convolution: given two real vectors $\ensuremath{\bm{\a}} = [a_0,\ensuremath{\ldots},a_{n-1}]^\ensuremath{\top}$ and $\ensuremath{\bm{\f}} = [f_0,\ensuremath{\ldots},f_{n-1}]^\ensuremath{\top}$, the \emph{cross-correlation} is defined by $\ensuremath{\bm{\a}} \ensuremath{\otimes} \ensuremath{\bm{\f}} = [c_0,\ensuremath{\ldots},c_{n-1}]^\ensuremath{\top}$ where
\[
c_k := \ensuremath{\sum}_{j=0}^{n-1} a_j f_{{\rm mod}(j+k,n)}.
\]
Cross-correlations can be interpreted as a heuristic means to identify when two vectors are close to being shifted versions of each other. Here we explore how the cross-correlation is equivalent to a convolution:

\rule{\textwidth}{1pt}
\textbf{Problem 3}  Explain why cross-correlation commutes with translations in the sense that $S(\ensuremath{\bm{\a}} \ensuremath{\otimes} \ensuremath{\bm{\f}} ) = a \ensuremath{\otimes} (S\ensuremath{\bm{\f}})$, which guarantees its equivalent to a convolution. For what $\ensuremath{\bm{\b}}$ do we have
\[
\ensuremath{\bm{\a}} \ensuremath{\otimes} \ensuremath{\bm{\f}}  = \ensuremath{\bm{\b}} \ensuremath{\star} \ensuremath{\bm{\f}}?
\]
\rule{\textwidth}{1pt}
We now move on to Chebyshev polynomials, which are a canonical example of orthogonal polynomials and allow us to use the approximation power of Fourier series for non-periodic functions. In the notes we investigated the standard Chebyshev polynomials $T_n(x) = \cos n \acos x$. Here we investigate a second kind of Chebyshev polynomials:
\[
U_n(x) := {\sin(n+1) \acos x \over \sqrt{1-x^2}}.
\]
These correspond to sine expansions in a similar way to how $T_n$ correspond to cosine expansions. We will show they are orthogonal polynomials by proving that they are (1) polynomials and (2) orthogonal with respect to a specific inner product.

We show first  that they are polynomials by proving they satisfy a three-term recurrence (a property satisfied by all orthogonal polynomials).

\rule{\textwidth}{1pt}
\textbf{Problem 4(a)} Show that $U_n$ are polynomials that satisfy
\begin{align*}
U_0(x) &= 1, \\
x U_0(x) &= U_1(x)/2, \\
x U_n(x) &= {U_{n-1}(x)+U_{n+1}(x) \over 2}, \qquad n > 0.
\end{align*}
Use this to deduce $U_5(x)$.

\rule{\textwidth}{1pt}
We can now show that this family of polynomials are orthogonal with respect to a specific weight, which is often called the \emph{semicircle weight}:

\rule{\textwidth}{1pt}
\textbf{Problem 4(b)} Prove that $U_n$ are orthogonal  with respect to
\[
\ensuremath{\int}_{-1}^1 f(x) g(x) \sqrt{1-x^2} \D x.
\]
\rule{\textwidth}{1pt}
Similar to the way that the derivative of a cosine is expressible as a sine, we can relate the derivative of $T_n$ to $U_n$, though where the degree decreases.  Here we prove this result (which is an example of a property satisfied by all families of \emph{classical orthogonal polynomials}):

\rule{\textwidth}{1pt}
\textbf{Problem 5}  Show for $n = 1, 2, \ensuremath{\ldots}$:
\[
T_n'(x) = n U_{n-1}(x).
\]
\rule{\textwidth}{1pt}
This connection can be used to differentiate functions, and as seen in the lab it can achieve much higher accuracy than divided difference techniques, even close to machine precision accuracy for nice enough functions.

This idea can be taken further and used to solve Ordinary Differential Equations (ODEs) to very high accuracy, in something called the \href{https://epubs.siam.org/doi/10.1137/120865458}{ultraspherical spectral method}.



\end{document}